\chapter{Complexity — Jacobi}
%%% generated by the blueprint pipeline from TCSlib.Complexity.ClassP.BinaryArith.Jacobi
\section{Overview}

Declarations of the Lean module \texttt{TCSlib.Complexity.ClassP.BinaryArith.Jacobi}.

\section{Declarations}

\begin{theorem}[Binary value modulo eight]
\lean{Complexity.BinArith.bitsVal_mod_eight}
\label{Complexity.BinArith.bitsVal_mod_eight}
\leanok
\uses{Complexity.BinArith.bitsVal_eq_headD, Complexity.LogProg.bitsVal}
\difficulty{4}
For every finite Boolean word $w$, the remainder of its binary value upon division by
$8$ is $\operatorname{val}(w_0)+2\operatorname{val}(w_1)+4\operatorname{val}(w_2)$,
where $w_i$ is the bit in position $i$ starting from the first position, and a missing
bit is taken to be false.
\end{theorem}

\begin{theorem}[Binary value modulo four]
\lean{Complexity.BinArith.mod_four_eq_three_iff}
\label{Complexity.BinArith.mod_four_eq_three_iff}
\leanok
\uses{Complexity.BinArith.bitsVal_mod_eight, Complexity.BinArith.bitsVal_mod_two, Complexity.LogProg.bitsVal}
\difficulty{5}
Let $w$ be a bit string whose binary value is odd. Its value is congruent to $3$ modulo
$4$ if and only if the second bit of $w$ is $1$, where a missing second bit is treated
as $0$.
\end{theorem}

\begin{definition}[Sign encoded by the first bit]
\lean{Complexity.BinArith.signOf}
\label{Complexity.BinArith.signOf}
\leanok
For a finite Boolean word $s$, define its sign to be $-1$ if its first bit is true, and
$1$ otherwise; an empty word is assigned sign $1$.
\end{definition}

\begin{definition}[Flip the first-bit sign]
\lean{Complexity.BinArith.flipSign}
\label{Complexity.BinArith.flipSign}
\leanok
For a finite Boolean word $s$, define its sign-flip word to consist of the single bit
opposite to its first bit; if $s$ is empty, use false as its first bit before flipping.
\end{definition}

\begin{theorem}[Negating a sign by flipping its bit]
\lean{Complexity.BinArith.signOf_flipSign}
\label{Complexity.BinArith.signOf_flipSign}
\leanok
\uses{Complexity.BinArith.flipSign, Complexity.BinArith.signOf}
\difficulty{3}
For any finite bit string $s$, assign sign $-1$ when its first bit is $1$ and sign $+1$
otherwise, treating an empty string as having first bit $0$. Replace that first bit by
its complement and discard all remaining bits; the sign of the resulting string is the
negative of the original sign.
\end{theorem}

\begin{definition}[Jacobi-symbol state update]
\lean{Complexity.BinArith.jacStep}
\label{Complexity.BinArith.jacStep}
\leanok
\uses{Complexity.BinArith.flipSign, Complexity.BinArith.mod, Complexity.LogProg.bitsVal, Complexity.pairFstD, Complexity.pairSndD, Turing.pairEncode}
For a Boolean word $S$, the update leaves $S$ unchanged when the value of the first
decoded component is zero. Otherwise, if the first bit of that component is true, it
pairs the remainder of that component modulo the encoded value of the second component
with a nested pair containing the original first component and the sign word, flipping
the sign word exactly when the second bits of the two components after the first are
both true. If the first bit is false, it drops that bit and pairs the resulting word
with a nested pair containing the first component of the second decoded component; in
this case it flips the sign word exactly when the second and third bits of that decoded
component differ.
\end{definition}

\begin{theorem}[Polynomial-Time Jacobi Step]
\lean{Complexity.BinArith.polyTimeComputable_jacStep}
\label{Complexity.BinArith.polyTimeComputable_jacStep}
\leanok
\uses{Complexity.BinArith.flipSign, Complexity.BinArith.jacStep, Complexity.BinArith.polyTimeComputable_drop, Complexity.BinArith.polyTimeComputable_isZero, Complexity.BinArith.polyTimeComputable_mod, Complexity.BinArith.polyTimeComputable_slot2, Complexity.BinArith.polyTimeComputable_slot3, Complexity.BinArith.polyTimeComputable_xor, Complexity.PolyTimeComputable, Complexity.PolyTimeComputable.comp, Complexity.PolyTimeComputable.pairEncode, Complexity.pairSndD, Complexity.polyTimeComputable_and, Complexity.polyTimeComputable_headD, Complexity.polyTimeComputable_id, Complexity.polyTimeComputable_ite, Complexity.polyTimeComputable_not, Complexity.polyTimeComputable_pairFstD, Complexity.polyTimeComputable_tail}
\difficulty{6}
The function on finite Boolean words given by the specified Jacobi-symbol iteration step
is polynomial-time computable.
\end{theorem}

\begin{definition}[Binary encoding of the Jacobi symbol]
\lean{Complexity.BinArith.jacobi}
\label{Complexity.BinArith.jacobi}
\leanok
\uses{Complexity.LogProg.bitsVal, Complexity.pairFstD, Complexity.pairSndD}
For a Boolean word $z$, take the Jacobi symbol of the values of the two components
decoded from $z$, add $1$, convert the result to a natural number, and represent that
natural number by its binary digits.
\end{definition}

\begin{definition}[Jacobi-symbol invariant]
\lean{Complexity.BinArith.JacInv}
\label{Complexity.BinArith.JacInv}
\leanok
\uses{Complexity.BinArith.signOf, Complexity.LogProg.bitsVal, Complexity.pairFstD, Complexity.pairSndD, Turing.pairEncode}
For Boolean words $z$ and $S$, the invariant holds if there exist Boolean words $A,N,s$
such that $S$ encodes $A$ paired with $N$ paired with $s$, the lengths of $A$ and $N$
are each at most the length of $z$, and $s$ has length at most one. Moreover, if the
value of the second component decoded from $z$ is odd, then the value of $N$ is odd and
the sign of $s$ times the Jacobi symbol of the values of $A$ and $N$ equals the Jacobi
symbol of the values of the two components decoded from $z$.
\end{definition}

\begin{definition}[Jacobi measure]
\lean{Complexity.BinArith.jacMeasure}
\label{Complexity.BinArith.jacMeasure}
\leanok
\uses{Complexity.LogProg.bitsVal, Complexity.pairFstD, Complexity.pairSndD}
For a Boolean list $S$, let $a$ be the natural number represented by its first decoded
component, and let $b$ be the natural number represented by the first decoded component
of the second decoded component of $S$. The measure is the binary size of $ab$, plus $1$
if $b\le a$ and $0$ otherwise.
\end{definition}

\begin{theorem}[Jacobi symbol of two modulo an odd integer]
\lean{Complexity.BinArith.jacobiSym_two_bits}
\label{Complexity.BinArith.jacobiSym_two_bits}
\leanok
\uses{Complexity.BinArith.bitsVal_mod_eight, Complexity.BinArith.bitsVal_mod_two, Complexity.LogProg.bitsVal}
\difficulty{5}
Let $N$ be a finite bit string whose represented integer is odd. The Jacobi symbol
$(2/N)$ is $-1$ if the bits in positions two and three of $N$ differ, and $1$ if they
agree; a missing bit is taken to be $0$.
\end{theorem}

\begin{theorem}[Jacobi symbol for an even numerator]
\lean{Complexity.BinArith.jacobiSym_even}
\label{Complexity.BinArith.jacobiSym_even}
\leanok
\difficulty{2}
For natural numbers $A,N$, if $A$ is even, then $(A/N)=(2/N)(A/2N)$.
\end{theorem}

\begin{theorem}[Jacobi symbol for an odd numerator]
\lean{Complexity.BinArith.jacobiSym_odd}
\label{Complexity.BinArith.jacobiSym_odd}
\leanok
\difficulty{3}
For natural numbers $A,N$, if both are odd, then $(A/N)=(-1)^{oldsymbol{1}_{A\equiv 3\
(\mathrm{mod}\ 4),\ N\equiv 3\ (\mathrm{mod}\ 4)}}(N\bmod A/A)$, where the prefactor is
$-1$ when both are congruent to $3$ modulo $4$, and $1$ otherwise.
\end{theorem}

\begin{theorem}[Preservation of the Jacobi invariant]
\lean{Complexity.BinArith.jacInv_step}
\label{Complexity.BinArith.jacInv_step}
\leanok
\uses{Complexity.BinArith.JacInv, Complexity.BinArith.bitsVal_drop_one, Complexity.BinArith.bitsVal_mod_two, Complexity.BinArith.flipSign, Complexity.BinArith.jacStep, Complexity.BinArith.jacobiSym_even, Complexity.BinArith.jacobiSym_odd, Complexity.BinArith.jacobiSym_two_bits, Complexity.BinArith.length_bits_le_of_le, Complexity.BinArith.mod, Complexity.BinArith.mod_four_eq_three_iff, Complexity.BinArith.signOf_flipSign, Complexity.LogProg.bitsVal, Complexity.LogProg.bitsVal_bits, Complexity.pairFstD_pairEncode, Complexity.pairSndD_pairEncode}
\difficulty{6}
Let $z$ and $S$ be bit strings satisfying the following condition: there are bit strings
$A,N,s$ with $S=\langle A,\langle N,s\rangle\rangle$, $|A|\le |z|$, $|N|\le |z|$, and
$|s|\le 1$, such that if the integer represented by the second projection of $z$ is odd,
then the integer represented by $N$ is odd and the product of the sign encoded by $s$
with the Jacobi symbol of the integers represented by $A,N$ equals the Jacobi symbol of
the integers represented by the first and second projections of $z$. After one Jacobi
step on $S$, the resulting state still satisfies this condition for the same $z$.
\end{theorem}

\begin{theorem}[Decrease of the Jacobi measure]
\lean{Complexity.BinArith.jacMeasure_step}
\label{Complexity.BinArith.jacMeasure_step}
\leanok
\uses{Complexity.BinArith.bitsVal_drop_one, Complexity.BinArith.bitsVal_mod_two, Complexity.BinArith.jacMeasure, Complexity.BinArith.jacStep, Complexity.BinArith.mod, Complexity.BinArith.size_lt_of_two_mul_le, Complexity.LogProg.bitsVal, Complexity.LogProg.bitsVal_bits, Complexity.pairFstD_pairEncode, Complexity.pairSndD_pairEncode, Turing.pairEncode}
\difficulty{5}
Let $A,N,s$ be finite bit strings, and suppose the integer represented by $N$ is odd. If
the Jacobi measure of the encoding of $A$ together with the encoding of $N$ and $s$ is
positive, one Jacobi step strictly decreases that measure. If the measure is zero, the
step leaves the encoding unchanged.
\end{theorem}

\begin{definition}[Jacobi output encoding]
\lean{Complexity.BinArith.jacOut}
\label{Complexity.BinArith.jacOut}
\leanok
\uses{Complexity.BinArith.beq, Complexity.pairFstD, Complexity.pairSndD, Turing.pairEncode}
For a Boolean list $S$, inspect the second component of its second decoded component. If
the encoding of the first component of that second decoded component, paired with the
one-element list $[\mathrm{true}]$, passes the Boolean test, return the empty list when
the inspected component defaults to false, and otherwise return
$[\mathrm{false},\mathrm{true}]$. If the test fails, return $[\mathrm{true}]$.
\end{definition}

\begin{definition}[Jacobi algorithm]
\lean{Complexity.BinArith.jacobiAlg}
\label{Complexity.BinArith.jacobiAlg}
\leanok
\uses{Complexity.BinArith.jacOut, Complexity.BinArith.jacStep, Complexity.pairFstD, Complexity.pairSndD, Turing.pairEncode}
For a Boolean list $z$, encode its first decoded component together with the encoding of
its second decoded component paired with $[\mathrm{false}]$. Apply the step function
$2(|z|+1)$ times to this encoding, where $|z|$ is the length of $z$, and then apply the
output function.
\end{definition}

\begin{theorem}[Polynomial-Time Jacobi Algorithm]
\lean{Complexity.BinArith.polyTimeComputable_jacobiAlg}
\label{Complexity.BinArith.polyTimeComputable_jacobiAlg}
\leanok
\uses{Complexity.BinArith.JacInv, Complexity.BinArith.iterate_invariant, Complexity.BinArith.jacInv_step, Complexity.BinArith.jacOut, Complexity.BinArith.jacobiAlg, Complexity.BinArith.polyTimeComputable_beq, Complexity.BinArith.polyTimeComputable_iterate, Complexity.BinArith.polyTimeComputable_jacStep, Complexity.BinArith.polyTimeComputable_slot2, Complexity.BinArith.polyTimeComputable_slot3, Complexity.BinArith.signOf, Complexity.PolyTimeComputable, Complexity.PolyTimeComputable.comp, Complexity.PolyTimeComputable.pairEncode, Complexity.length_pairFstD_le, Complexity.length_pairSndD_le, Complexity.pairFstD, Complexity.pairSndD, Complexity.polyTimeComputable_const, Complexity.polyTimeComputable_headD, Complexity.polyTimeComputable_ite, Complexity.polyTimeComputable_pairFstD, Complexity.polyTimeComputable_pairSndD, Turing.length_pairEncode, Turing.pairEncode}
\difficulty{4}
The function on finite Boolean words given by the specified Jacobi algorithm is
polynomial-time computable.
\end{theorem}

\begin{theorem}[Correctness of the Jacobi algorithm]
\lean{Complexity.BinArith.jacobiAlg_eq}
\label{Complexity.BinArith.jacobiAlg_eq}
\leanok
\uses{Complexity.BinArith.JacInv, Complexity.BinArith.beq, Complexity.BinArith.iterate_invariant, Complexity.BinArith.iterate_measure_eq_zero, Complexity.BinArith.jacInv_step, Complexity.BinArith.jacMeasure, Complexity.BinArith.jacMeasure_step, Complexity.BinArith.jacOut, Complexity.BinArith.jacStep, Complexity.BinArith.jacobi, Complexity.BinArith.jacobiAlg, Complexity.BinArith.signOf, Complexity.LogProg.bitsVal, Complexity.LogProg.bitsVal_lt_two_pow, Complexity.length_pairFstD_le, Complexity.length_pairSndD_le, Complexity.length_pair_components_le, Complexity.pairFstD, Complexity.pairFstD_pairEncode, Complexity.pairSndD, Complexity.pairSndD_pairEncode, Turing.pairEncode}
\difficulty{6}
Let $z$ be a finite bit string whose second decoded component represents an odd integer.
The output of the Jacobi algorithm on $z$ is the binary representation of one plus the
Jacobi symbol of the integers represented by the first and second decoded components of
$z$.
\end{theorem}

\begin{theorem}[Jacobi algorithm on encoded integers]
\lean{Complexity.BinArith.jacobiAlg_pairEncode}
\label{Complexity.BinArith.jacobiAlg_pairEncode}
\leanok
\uses{Complexity.BinArith.jacobi, Complexity.BinArith.jacobiAlg, Complexity.BinArith.jacobiAlg_eq, Complexity.LogProg.bitsVal_bits, Turing.pairEncode}
\difficulty{1}
For natural numbers $a$ and $n$ with $n$ odd, applying the Jacobi algorithm to the pair
encoding of their binary expansions yields the binary expansion of
$(\operatorname{jacobiSym}(a,n)+1).toNat$.
\end{theorem}

\begin{definition}[Jacobi modular algorithm]
\lean{Complexity.BinArith.jacobiModAlg}
\label{Complexity.BinArith.jacobiModAlg}
\leanok
\uses{Complexity.BinArith.add, Complexity.BinArith.jacobiAlg, Complexity.BinArith.mod, Complexity.BinArith.sub, Complexity.pairSndD, Turing.pairEncode}
For a Boolean list $z$, apply the modular operation to the encoding of the subtraction
of $[\mathrm{true}]$ from the encoding of the addition of the Jacobi algorithm's output
on $z$ and the second decoded component of $z$, paired with that same second decoded
component.
\end{definition}

\begin{theorem}[Polynomial-Time Jacobi Modular Algorithm]
\lean{Complexity.BinArith.polyTimeComputable_jacobiModAlg}
\label{Complexity.BinArith.polyTimeComputable_jacobiModAlg}
\leanok
\uses{Complexity.BinArith.jacobiModAlg, Complexity.BinArith.polyTimeComputable_add, Complexity.BinArith.polyTimeComputable_jacobiAlg, Complexity.BinArith.polyTimeComputable_mod, Complexity.BinArith.polyTimeComputable_sub, Complexity.PolyTimeComputable, Complexity.PolyTimeComputable.comp, Complexity.PolyTimeComputable.pairEncode, Complexity.polyTimeComputable_const, Complexity.polyTimeComputable_pairSndD}
\difficulty{1}
The function on finite Boolean words given by the specified Jacobi modular algorithm is
polynomial-time computable.
\end{theorem}

\begin{theorem}[Jacobi symbol modulo an odd integer]
\lean{Complexity.BinArith.jacobiModAlg_pairEncode}
\label{Complexity.BinArith.jacobiModAlg_pairEncode}
\leanok
\uses{Complexity.BinArith.add_pairEncode, Complexity.BinArith.jacobiAlg_pairEncode, Complexity.BinArith.jacobiModAlg, Complexity.BinArith.mod, Complexity.BinArith.sub, Complexity.LogProg.bitsVal, Complexity.LogProg.bitsVal_bits, Complexity.pairFstD_pairEncode, Complexity.pairSndD_pairEncode, Turing.pairEncode}
\difficulty{5}
For natural numbers $a$ and $n$ with $n$ odd, applying the modular Jacobi algorithm to
the pairing of their binary expansions yields the binary expansion of the least
nonnegative residue of the Jacobi symbol $(a/n)$ modulo $n$.
\end{theorem}
